\documentclass[11pt,a4paper]{article}

% --- Encoding ---
\usepackage[utf8]{inputenc}
\usepackage[T1]{fontenc}
\usepackage[spanish,es-noshorthands]{babel}
\usepackage{lmodern}

% --- Layout ---
\usepackage[margin=2.5cm]{geometry}
\usepackage{parskip}
\setlength{\parindent}{0pt}
\setlength{\parskip}{0.5em plus 0.2em minus 0.1em}

% --- Tables and graphics ---
\usepackage{booktabs}
\usepackage{array}
\usepackage{enumitem}
\usepackage{xcolor}
\usepackage{graphicx}
\usepackage{tikz}
\usetikzlibrary{positioning,shapes.geometric,arrows.meta,fit,backgrounds,calc,decorations.pathreplacing}
\usepackage{caption}
\usepackage{subcaption}
\usepackage{amsmath,amssymb}

% --- Soporte para caracteres Unicode usados en formulas en texto plano ---
\usepackage{newunicodechar}
\newunicodechar{→}{\ensuremath{\rightarrow}}
\newunicodechar{≥}{\ensuremath{\geq}}
\newunicodechar{≤}{\ensuremath{\leq}}
\newunicodechar{×}{\ensuremath{\times}}
\newunicodechar{÷}{\ensuremath{\div}}
\newunicodechar{≈}{\ensuremath{\approx}}
\newunicodechar{⌈}{\ensuremath{\lceil}}
\newunicodechar{⌉}{\ensuremath{\rceil}}
\newunicodechar{✓}{\ensuremath{\checkmark}}
\newunicodechar{∀}{\ensuremath{\forall}}
\newunicodechar{∧}{\ensuremath{\wedge}}

% --- Code listings ---
\usepackage{listings}
\definecolor{codegray}{rgb}{0.96,0.96,0.96}
\definecolor{codeborder}{rgb}{0.85,0.85,0.85}
\definecolor{codecomment}{rgb}{0.30,0.55,0.30}
\definecolor{codekeyword}{rgb}{0.10,0.30,0.70}
\definecolor{codestring}{rgb}{0.65,0.15,0.10}

\lstdefinestyle{cstyle}{
    language=C,
    backgroundcolor=\color{codegray},
    basicstyle=\ttfamily\footnotesize,
    breaklines=true,
    frame=single,
    rulecolor=\color{codeborder},
    showstringspaces=false,
    keywordstyle=\color{codekeyword}\bfseries,
    commentstyle=\color{codecomment}\itshape,
    stringstyle=\color{codestring},
    aboveskip=1em,
    belowskip=1em,
}

\lstdefinestyle{shellstyle}{
    language=bash,
    backgroundcolor=\color{codegray},
    basicstyle=\ttfamily\small,
    breaklines=true,
    frame=single,
    rulecolor=\color{codeborder},
    showstringspaces=false,
    keywordstyle=\color{codekeyword}\bfseries,
    commentstyle=\color{codecomment}\itshape,
    stringstyle=\color{codestring},
    aboveskip=1em,
    belowskip=1em,
    columns=fullflexible,
    upquote=true,
}

\lstdefinestyle{textstyle}{
    backgroundcolor=\color{codegray},
    basicstyle=\ttfamily\footnotesize,
    breaklines=true,
    frame=single,
    rulecolor=\color{codeborder},
    showstringspaces=false,
    aboveskip=1em,
    belowskip=1em,
    columns=fullflexible,
}

\lstset{style=textstyle}

% --- Definition / Idea / Important boxes ---
\usepackage[most]{tcolorbox}
\newtcolorbox{idea}[1][]{
    colback=orange!5!white,
    colframe=orange!60!black,
    fonttitle=\bfseries,
    title=Idea intuitiva,
    sharp corners=southwest,
    rounded corners=northwest,
    arc=2pt,
    boxrule=0.6pt,
    left=8pt,right=8pt,top=4pt,bottom=4pt,
    #1
}
\newtcolorbox{importante}[1][]{
    colback=red!5!white,
    colframe=red!50!black,
    fonttitle=\bfseries,
    title=Importante,
    sharp corners=southwest,
    rounded corners=northwest,
    arc=2pt,
    boxrule=0.6pt,
    left=8pt,right=8pt,top=4pt,bottom=4pt,
    #1
}
\newtcolorbox{respuesta}[1][]{
    colback=green!5!white,
    colframe=green!40!black,
    fonttitle=\bfseries,
    title=Respuesta,
    sharp corners=southwest,
    rounded corners=northwest,
    arc=2pt,
    boxrule=0.6pt,
    left=8pt,right=8pt,top=4pt,bottom=4pt,
    #1
}

% --- Hyperlinks ---
\usepackage[colorlinks=true,
            linkcolor=blue!60!black,
            urlcolor=blue!50!black,
            citecolor=blue!60!black]{hyperref}

% --- Color palette ---
\definecolor{c1}{HTML}{4C72B0}
\definecolor{c2}{HTML}{DD8452}
\definecolor{c3}{HTML}{55A868}
\definecolor{c4}{HTML}{8172B2}
\definecolor{cred}{HTML}{C44E52}
\definecolor{cfree}{HTML}{8C8C8C}

% --- Title block ---
\title{
    \vspace{-2.5em}
    \textbf{Parcial 2 --- 2026} \\[0.3em]
    \large Sistemas Operativos / Sistemas Operativos y Redes
}
\author{
    Tomás Rodeghiero \\
    \small UNRC --- Universidad Nacional de Río Cuarto
}
\date{2026}

\begin{document}

\maketitle

\begin{abstract}
\noindent
Resolución completa del Parcial 2 del año 2026 de la cátedra de
Sistemas Operativos (Lic.\ en Ciencias de la Computación, UNRC). El
parcial consta de cinco ejercicios que cubren: (i) verdadero/falso
sobre gestión de heap, fragmentación, \emph{best fit}, \texttt{mmap}
y paginación; (ii) evolución de un \emph{buddy system} con una
secuencia de \texttt{malloc}/\texttt{free}; (iii) diseño de tabla de
páginas para un proceso, traducción de direcciones e identificación de
una página de guarda para detectar \emph{stack overflow}; (iv)
representación de un sistema de archivos FAT32 dado el árbol de un
volumen pequeño con un archivo binario de 2 bloques y un directorio
vacío; (v) verdadero/falso sobre privilegios de usuario, \texttt{sudo},
contraseñas y permisos UGO. Cada ejercicio se desarrolla con la
estructura \textbf{Enunciado $\to$ Análisis $\to$ Desarrollo $\to$
Respuesta}.
\end{abstract}

\tableofcontents
\bigskip

% =====================================================================
\section{Ejercicio 1 --- V/F sobre gestión de memoria (1.5 pts)}
% =====================================================================

\textbf{Enunciado.} Determinar la verdad o falsedad de las siguientes
sentencias. Justificar las falsas en forma breve y con claridad.

\subsection*{1.a --- Heap de bloques variables y fragmentación}

\emph{``El gestor de un heap de bloques de tamaño variable tiene
problemas de fragmentación.''}

\begin{respuesta}
\textbf{V}. Un gestor de heap con bloques de tamaño variable (como el
\texttt{malloc} clásico, \emph{dlmalloc}, \emph{ptmalloc}) sufre
inherentemente \textbf{fragmentación externa}: con el correr de
muchas asignaciones y liberaciones de tamaños distintos, los huecos
libres se fragmentan en bloques pequeños no contiguos, y puede pasar
que la suma total de memoria libre sea grande pero no haya un único
hueco contiguo suficiente para satisfacer un pedido.

También puede haber \textbf{fragmentación interna} si el gestor
redondea cada pedido hacia arriba (a un múltiplo del tamaño mínimo).
\end{respuesta}

\subsection*{1.b --- Bitset para bloques de tamaño fijo}

\emph{``Un gestor de bloques de tamaño fijo puede manejar los bloques
libres/ocupados con un bitset.''}

\begin{respuesta}
\textbf{V}. Cuando todos los bloques tienen el mismo tamaño, alcanza
con \textbf{un bit por bloque} (0 = libre, 1 = ocupado). El
\emph{bitmap} (o \emph{bitset}) es el método clásico de
\emph{slab allocators}, asignadores de \emph{frames} de memoria
física y administradores de \emph{inodes}/bloques de disco.
\end{respuesta}

\subsection*{1.c --- \emph{Best fit} y lista de libres ordenada}

\emph{``La estrategia \emph{best fit} en un \texttt{malloc(s)} asigna
un bloque libre $b$ con \texttt{size(b)} ≥ $s$ y $∀$ $b'$:
\texttt{free(b')} $∧$ \texttt{size(b')} >= \texttt{size(b)}. Además
conviene mantener la lista de bloques libres ordenadas de menor a
mayor tamaño.''}

\begin{respuesta}
\textbf{V}. \emph{Best fit} busca, entre los bloques libres que
satisfacen el pedido (\texttt{size(b) ≥ s}), el de \textbf{menor
tamaño} (es decir, el que mejor ``calza''). La intuición es que esto
deja los bloques grandes intactos para futuros pedidos grandes,
reduciendo la probabilidad de fragmentar uno innecesariamente.

Mantener la lista de libres ordenada \emph{de menor a mayor tamaño}
es la implementación natural: se recorre la lista hasta encontrar el
primer bloque $\geq s$ y ese es el \emph{best fit}, en tiempo
proporcional a la cantidad de bloques de tamaño menor que $s$. (En la
práctica se usan estructuras más eficientes, como árboles balanceados
ordenados por tamaño, para hacerlo en $O(\log n)$.)
\end{respuesta}

\subsection*{1.d --- \texttt{mmap} carga todo el archivo a memoria}

\emph{``La llamada al sistema \texttt{mmap(file, ...)} carga todo el
archivo en memoria y retorna la dirección lógica del comienzo del
área de memoria asignada.''}

\begin{respuesta}
\textbf{F}. \texttt{mmap} \textbf{no carga el archivo entero} en
memoria física al momento de la llamada. Lo que hace es crear un
\emph{mapeo} en el espacio de direcciones virtuales del proceso: se
agregan entradas en la tabla de páginas marcadas como
\emph{no presentes} ($P=0$) que asocian rangos virtuales con bloques
del archivo en disco. Las páginas se traen a memoria \textbf{bajo
demanda} (\emph{demand paging}): cuando el proceso accede a una
dirección dentro del mapeo, el hardware genera un \emph{page fault};
el kernel atiende la falla leyendo el bloque correspondiente del
archivo y poblando la página.

Sí es cierto que \texttt{mmap} devuelve la dirección lógica del
comienzo del área mapeada (eso está bien). El problema de la
afirmación es la palabra ``carga todo el archivo'': la carga es
\emph{perezosa}, no inmediata.
\end{respuesta}

\subsection*{1.e --- Ventaja de la paginación}

\emph{``La paginación tiene la ventaja que los bloques de un proceso
en memoria (código, datos globales, stack, etc) no requieren estar
físicamente contiguos.''}

\begin{respuesta}
\textbf{V}. Ésta es \textbf{la} ventaja característica de la
paginación. El espacio de direcciones virtual del proceso es
\emph{lógicamente contiguo}, pero la MMU traduce cada página virtual
a una página física independiente: las páginas físicas de un mismo
proceso pueden estar dispersas en cualquier parte de la RAM. Esto
elimina la fragmentación externa que sufren las particiones
variables y permite asignar memoria a cualquier proceso siempre que
haya \emph{frames} libres, sin importar su distribución espacial.
\end{respuesta}

% =====================================================================
\section{Ejercicio 2 --- Evolución de un \emph{buddy system} (2 pts)}
% =====================================================================

\textbf{Enunciado.} Mostrar la evolución de los bloques ocupados y
libres de un \emph{buddy system} con bloques de tamaño entre $2^6$ y
$2^9$. El tamaño del heap es de $2$\,KB. Secuencia de operaciones:

\begin{enumerate}[leftmargin=*,label=\Roman*.,nosep]
    \item \texttt{ptr1 = malloc(500)}
    \item \texttt{ptr2 = malloc(8)}
    \item \texttt{ptr3 = malloc(120)}
    \item \texttt{free(ptr2)}
    \item \texttt{ptr4 = alloc(600)}
\end{enumerate}

\subsection*{Interpretación del enunciado}

El heap mide $2$\,KB $= 2^{11}$ bytes y los tamaños de bloque
mencionados son $2^6 = 64$ y $2^9 = 512$. Tomamos esto como una
indicación de la \emph{granularidad típica} de las asignaciones: el
heap se inicializa como un único bloque grande de $2^{11} = 2048$
bytes y el \emph{buddy system} puede subdividirlo recursivamente
hasta alcanzar el mínimo de $64$ bytes, o consolidar bloques
contiguos hasta el máximo de $2048$.

\begin{idea}
Bajo esta interpretación, el último pedido \texttt{alloc(600)} se
satisface con un bloque de $1024 = 2^{10}$ bytes y el ejercicio
muestra toda la mecánica del buddy. Si en cambio se interpretara
estrictamente que el bloque máximo es $2^9 = 512$, el último pedido
fallaría por no haber ningún bloque de $1024$ disponible; la versión
adoptada acá es la pedagógicamente más rica.
\end{idea}

\subsection*{Algoritmo del \emph{buddy system}}

\begin{itemize}[leftmargin=*,nosep]
    \item Para un pedido de tamaño $s$, se busca el menor bloque
        libre de tamaño $2^k$ con $2^k \geq s$.
    \item Si el bloque encontrado es \emph{demasiado grande}, se
        divide a la mitad recursivamente hasta llegar al tamaño justo.
    \item Al liberar un bloque, se verifica si su \emph{buddy}
        (el bloque vecino del mismo tamaño con el bit complementario
        de offset) está libre; si lo está, se fusionan en un bloque
        del doble del tamaño. La fusión se repite mientras sea posible.
\end{itemize}

\subsection*{Estado inicial}

Heap libre de $2048$ bytes en el offset $0$.

\begin{center}
\begin{tikzpicture}[xscale=0.13,yscale=0.6,every node/.style={font=\ttfamily\small}]
\draw[fill=cfree!30] (0,0) rectangle (50,1) node[midway,font=\bfseries\small]{2048 (libre)};
\draw (0,0) -- (0,-0.3) node[below,font=\scriptsize]{0};
\draw (50,0) -- (50,-0.3) node[below,font=\scriptsize]{2048};
\end{tikzpicture}
\end{center}

\subsection*{Paso I: \texttt{ptr1 = malloc(500)}}

$500$ bytes redondeado hacia la siguiente potencia de $2$ que sea
\(\geq 500\): es $2^9 = 512$.

Necesitamos un bloque libre de $512$. El único bloque libre es de
$2048$. Lo dividimos sucesivamente:

\begin{lstlisting}[style=textstyle]
2048 -> 1024 + 1024     (tomo el primer 1024)
1024 -> 512 + 512       (tomo el primer 512)
\end{lstlisting}

Asigno el bloque de $512$ en offset $0$ a \texttt{ptr1}.

\begin{center}
\begin{tikzpicture}[xscale=0.13,yscale=0.6,every node/.style={font=\ttfamily\scriptsize}]
\draw[fill=c1!50]    (0,0)    rectangle (12.5,1) node[midway,white,font=\bfseries]{ptr1: 512};
\draw[fill=cfree!30] (12.5,0) rectangle (25,1)   node[midway,font=\bfseries]{512 libre};
\draw[fill=cfree!30] (25,0)   rectangle (50,1)   node[midway,font=\bfseries]{1024 libre};
\foreach \x in {0,12.5,25,50}
    \draw (\x,0) -- (\x,-0.3) node[below]{\pgfmathprintnumber{\x}};
\end{tikzpicture}
\end{center}

\subsection*{Paso II: \texttt{ptr2 = malloc(8)}}

$8$ bytes redondeado al mínimo permitido $2^6 = 64$.

El bloque libre más chico disponible es $512$ (en offset $512$).
Subdividimos hasta llegar a $64$:

\begin{lstlisting}[style=textstyle]
512 -> 256 + 256        (tomo el primer 256, offset 512)
256 -> 128 + 128        (tomo el primer 128, offset 512)
128 -> 64 + 64          (tomo el primer 64, offset 512)
\end{lstlisting}

\texttt{ptr2} apunta al bloque de $64$ en offset $512$.

\begin{center}
\begin{tikzpicture}[xscale=0.13,yscale=0.6,every node/.style={font=\ttfamily\scriptsize}]
\draw[fill=c1!50]    (0,0)     rectangle (12.5,1)   node[midway,white,font=\bfseries]{ptr1: 512};
\draw[fill=c2!50]    (12.5,0)  rectangle (14,1)     node[midway,white,font=\bfseries]{p2: 64};
\draw[fill=cfree!30] (14,0)    rectangle (15.5,1)   node[midway]{64};
\draw[fill=cfree!30] (15.5,0)  rectangle (18.625,1) node[midway]{128};
\draw[fill=cfree!30] (18.625,0) rectangle (25,1)    node[midway]{256};
\draw[fill=cfree!30] (25,0)    rectangle (50,1)     node[midway,font=\bfseries]{1024};
\draw (0,0)     -- (0,-0.3)     node[below]{0};
\draw (12.5,0)  -- (12.5,-0.3)  node[below]{512};
\draw (14,0)    -- (14,-0.3)    node[below]{576};
\draw (15.5,0)  -- (15.5,-0.3)  node[below]{640};
\draw (18.625,0)-- (18.625,-0.3)node[below]{768};
\draw (25,0)    -- (25,-0.3)    node[below]{1024};
\draw (50,0)    -- (50,-0.3)    node[below]{2048};
\end{tikzpicture}
\end{center}

\subsection*{Paso III: \texttt{ptr3 = malloc(120)}}

$120$ se redondea a $2^7 = 128$. Hay un bloque libre de $128$ en
offset $640$ que coincide exactamente. Lo asigno a \texttt{ptr3}, sin
necesidad de subdividir.

\begin{center}
\begin{tikzpicture}[xscale=0.13,yscale=0.6,every node/.style={font=\ttfamily\scriptsize}]
\draw[fill=c1!50]    (0,0)     rectangle (12.5,1)   node[midway,white,font=\bfseries]{ptr1: 512};
\draw[fill=c2!50]    (12.5,0)  rectangle (14,1)     node[midway,white,font=\bfseries]{p2: 64};
\draw[fill=cfree!30] (14,0)    rectangle (15.5,1)   node[midway]{64};
\draw[fill=c3!50]    (15.5,0)  rectangle (18.625,1) node[midway,white,font=\bfseries]{p3: 128};
\draw[fill=cfree!30] (18.625,0) rectangle (25,1)    node[midway]{256};
\draw[fill=cfree!30] (25,0)    rectangle (50,1)     node[midway,font=\bfseries]{1024};
\draw (0,0)     -- (0,-0.3)     node[below]{0};
\draw (12.5,0)  -- (12.5,-0.3)  node[below]{512};
\draw (14,0)    -- (14,-0.3)    node[below]{576};
\draw (15.5,0)  -- (15.5,-0.3)  node[below]{640};
\draw (18.625,0)-- (18.625,-0.3)node[below]{768};
\draw (25,0)    -- (25,-0.3)    node[below]{1024};
\draw (50,0)    -- (50,-0.3)    node[below]{2048};
\end{tikzpicture}
\end{center}

\subsection*{Paso IV: \texttt{free(ptr2)}}

Libero el bloque de $64$ en offset $512$. El buddy de un bloque de
tamaño $64$ en offset $512$ está en offset $512 \oplus 64 = 576$.
El bloque en $576$ (de $64$ bytes) está libre $\to$ \textbf{fusiono}:

\quad \texttt{libre 64@512 + libre 64@576 → libre 128@512}

Intento volver a fusionar el nuevo bloque de $128$ en offset $512$.
Su buddy está en offset $512 \oplus 128 = 640$. Pero en $640$ está
\texttt{ptr3} (ocupado). \textbf{No se puede fusionar más.}

\begin{center}
\begin{tikzpicture}[xscale=0.13,yscale=0.6,every node/.style={font=\ttfamily\scriptsize}]
\draw[fill=c1!50]    (0,0)     rectangle (12.5,1)   node[midway,white,font=\bfseries]{ptr1: 512};
\draw[fill=cfree!30] (12.5,0)  rectangle (15.5,1)   node[midway,font=\bfseries]{128 libre};
\draw[fill=c3!50]    (15.5,0)  rectangle (18.625,1) node[midway,white,font=\bfseries]{p3: 128};
\draw[fill=cfree!30] (18.625,0) rectangle (25,1)    node[midway]{256};
\draw[fill=cfree!30] (25,0)    rectangle (50,1)     node[midway,font=\bfseries]{1024};
\draw (0,0)     -- (0,-0.3)     node[below]{0};
\draw (12.5,0)  -- (12.5,-0.3)  node[below]{512};
\draw (15.5,0)  -- (15.5,-0.3)  node[below]{640};
\draw (18.625,0)-- (18.625,-0.3)node[below]{768};
\draw (25,0)    -- (25,-0.3)    node[below]{1024};
\draw (50,0)    -- (50,-0.3)    node[below]{2048};
\end{tikzpicture}
\end{center}

\subsection*{Paso V: \texttt{ptr4 = alloc(600)}}

$600$ se redondea a $2^{10} = 1024$. Hay un bloque libre de $1024$
exactamente en offset $1024$. Lo asigno a \texttt{ptr4}.

\begin{center}
\begin{tikzpicture}[xscale=0.13,yscale=0.6,every node/.style={font=\ttfamily\scriptsize}]
\draw[fill=c1!50]    (0,0)     rectangle (12.5,1)   node[midway,white,font=\bfseries]{ptr1: 512};
\draw[fill=cfree!30] (12.5,0)  rectangle (15.5,1)   node[midway,font=\bfseries]{128 libre};
\draw[fill=c3!50]    (15.5,0)  rectangle (18.625,1) node[midway,white,font=\bfseries]{p3: 128};
\draw[fill=cfree!30] (18.625,0) rectangle (25,1)    node[midway]{256};
\draw[fill=c4!50]    (25,0)    rectangle (50,1)     node[midway,white,font=\bfseries]{ptr4: 1024};
\draw (0,0)     -- (0,-0.3)     node[below]{0};
\draw (12.5,0)  -- (12.5,-0.3)  node[below]{512};
\draw (15.5,0)  -- (15.5,-0.3)  node[below]{640};
\draw (18.625,0)-- (18.625,-0.3)node[below]{768};
\draw (25,0)    -- (25,-0.3)    node[below]{1024};
\draw (50,0)    -- (50,-0.3)    node[below]{2048};
\end{tikzpicture}
\end{center}

\subsection*{Estado final --- resumen}

\begin{center}
\small
\begin{tabular}{@{}cccl@{}}
\toprule
Offset & Tamaño & Estado & Apunta a / nota \\
\midrule
$0$    & $512$  & ocupado & \texttt{ptr1} (pedido original: 500 B; frag.\ interna 12 B) \\
$512$  & $128$  & libre   & --- (proviene de la fusión del paso IV) \\
$640$  & $128$  & ocupado & \texttt{ptr3} (pedido: 120 B; frag.\ interna 8 B) \\
$768$  & $256$  & libre   & --- \\
$1024$ & $1024$ & ocupado & \texttt{ptr4} (pedido: 600 B; frag.\ interna 424 B) \\
\bottomrule
\end{tabular}
\end{center}

\begin{respuesta}
\textbf{Punteros activos al final:}

\begin{itemize}[leftmargin=*,nosep]
    \item \texttt{ptr1} $\to$ bloque de $512$ bytes en offset $0$.
    \item \texttt{ptr2} $\to$ liberado (no apunta a nada válido).
    \item \texttt{ptr3} $\to$ bloque de $128$ bytes en offset $640$.
    \item \texttt{ptr4} $\to$ bloque de $1024$ bytes en offset $1024$.
\end{itemize}

\textbf{Bloques libres:} $128$ B en offset $512$ y $256$ B en offset
$768$. Memoria libre total: $384$ bytes (pero \emph{ningún bloque
contiguo} mayor a $256$).
\end{respuesta}

% =====================================================================
\section{Ejercicio 3 --- Paginación con páginas de 2 KB (3 pts)}
% =====================================================================

\textbf{Enunciado.} CPU de 32 bits con espacio de direcciones de
$1$\,MB $= 2^{20}$ y páginas de $2$\,KB.

\subsection*{Análisis de los parámetros básicos}

\begin{itemize}[leftmargin=*,nosep]
    \item Tamaño de página: $2$\,KB $= 2^{11}$ bytes.
    \item Bits para el \texttt{offset} dentro de la página: $11$.
    \item Bits para el \texttt{vpn} (\emph{virtual page number}):
        $20 - 11 = 9$.
    \item Cantidad de páginas en el espacio total: $2^{20} / 2^{11}
        = 2^9 = 512$.
\end{itemize}

\subsection*{3.I --- Formato de la tabla de páginas}

\begin{itemize}[leftmargin=*,nosep]
    \item \textbf{Cantidad de entradas}: $2^9 = 512$.
    \item Cada entrada (\texttt{pte}) almacena:
    \begin{itemize}[nosep]
        \item \texttt{ppn} (\emph{physical page number}): si el
            espacio físico también es de $1$\,MB, hacen falta $9$
            bits.
        \item \texttt{flags}: $V$ (válida), $P$ (presente en RAM),
            $W$ (escribible), $X$ (ejecutable). Son $4$ bits.
    \end{itemize}
    \item Total por entrada: $9 + 4 = 13$ bits útiles; redondeado a
        \textbf{$16$ bits ($2$ bytes)} por motivos de alineación.
    \item \textbf{Tamaño de la tabla}: $512 \times 2$ B $= 1024$ B $=
        1$\,KB. Cabe holgadamente en una página.
\end{itemize}

\begin{center}
\begin{tikzpicture}[every node/.style={font=\ttfamily\small},
    cell/.style={draw,minimum height=0.7cm,align=center}]
\node[cell,minimum width=2.4cm,fill=c1!20] (ppn) at (0,0) {ppn};
\node[cell,minimum width=0.6cm,fill=c3!30,right=0pt of ppn] (V) {V};
\node[cell,minimum width=0.6cm,fill=c3!30,right=0pt of V]   (P) {P};
\node[cell,minimum width=0.6cm,fill=c3!30,right=0pt of P]   (W) {W};
\node[cell,minimum width=0.6cm,fill=c3!30,right=0pt of W]   (X) {X};
\node[cell,minimum width=1.0cm,right=0pt of X]              (res) {res.};
\node[font=\scriptsize] at ($(ppn.north)+(0,0.25)$) {9 bits};
\node[font=\scriptsize] at ($(V.north)+(0,0.25)$) {1};
\node[font=\scriptsize] at ($(P.north)+(0,0.25)$) {1};
\node[font=\scriptsize] at ($(W.north)+(0,0.25)$) {1};
\node[font=\scriptsize] at ($(X.north)+(0,0.25)$) {1};
\node[font=\scriptsize] at ($(res.north)+(0,0.25)$) {3};
\end{tikzpicture}
\end{center}

\subsection*{3.II --- Interpretación de una dirección virtual}

Las direcciones virtuales tienen $20$ bits útiles. La MMU las parte
en dos campos:

\begin{center}
\begin{tikzpicture}[every node/.style={font=\ttfamily\small},
    cell/.style={draw,minimum height=0.7cm,align=center}]
\node[cell,minimum width=2.6cm,fill=c1!20] (vpn) at (0,0) {vpn};
\node[cell,minimum width=4.6cm,fill=c2!20,right=0pt of vpn] (off) {offset};
\node[font=\scriptsize] at ($(vpn.north)+(0,0.25)$) {9 bits};
\node[font=\scriptsize] at ($(off.north)+(0,0.25)$) {11 bits};
\node[font=\scriptsize] at ($(vpn.north west)+(0.1,-0.05)$) {19};
\node[font=\scriptsize] at ($(vpn.north east)+(-0.15,-0.05)$) {11};
\node[font=\scriptsize] at ($(off.north west)+(0.1,-0.05)$) {10};
\node[font=\scriptsize] at ($(off.north east)+(-0.05,-0.05)$) {0};
\end{tikzpicture}
\end{center}

\begin{itemize}[leftmargin=*,nosep]
    \item Los \textbf{9 bits altos} forman el \texttt{vpn}: indexan
        una de las $512$ entradas de la tabla.
    \item Los \textbf{11 bits bajos} son el \texttt{offset} dentro de
        la página de $2$\,KB.
\end{itemize}

\subsection*{3.III --- Tabla de páginas de un proceso concreto}

\textbf{Datos del proceso.}
\begin{itemize}[leftmargin=*,nosep]
    \item Espacio de direcciones: $[0, 512\text{\,KB})$. Esto son
        $512\text{\,KB} \div 2\text{\,KB} = 256$ páginas (\texttt{vpn}
        $0$ a $255$).
    \item $2$ páginas de código $\to$ vpn $0$ y $1$.
    \item $1$ página de datos globales $\to$ vpn $2$.
    \item $2$ páginas de stack \emph{al final del espacio}
        $\to$ vpn $254$ y $255$.
\end{itemize}

\textbf{Tabla.} Los \texttt{ppn} concretos los asigna el SO según los
frames físicos libres; usamos valores arbitrarios como ejemplo.

\begin{center}
\small
\begin{tabular}{@{}cccccccl@{}}
\toprule
\texttt{vpn} & \texttt{ppn} & V & P & W & X & U & Área \\
\midrule
$0$     & \texttt{0x010} & 1 & 1 & 0 & 1 & 1 & código (página 1) \\
$1$     & \texttt{0x011} & 1 & 1 & 0 & 1 & 1 & código (página 2) \\
$2$     & \texttt{0x020} & 1 & 1 & 1 & 0 & 1 & datos globales \\
$3 \ldots 253$ & ---       & 0 & 0 & 0 & 0 & 0 & inválido (gap) \\
$254$   & \texttt{0x030} & 1 & 1 & 1 & 0 & 1 & stack (página 1) \\
$255$   & \texttt{0x031} & 1 & 1 & 1 & 0 & 1 & stack (página 2) \\
\bottomrule
\end{tabular}
\end{center}

\textbf{Justificación de las flags:}

\begin{itemize}[leftmargin=*,nosep]
    \item \textbf{Código}: $V=P=1$, $W=0$ (no se modifica),
        $X=1$ (ejecutable).
    \item \textbf{Datos / stack}: $V=P=1$, $W=1$ (escribible),
        $X=0$ (no se ejecuta como código; protección contra
        explotación).
    \item \textbf{Gap (3--253)}: $V=0$. Cualquier acceso a una de
        esas direcciones genera un \emph{page fault}.
\end{itemize}

\subsection*{3.IV --- ¿A qué página física se mapea \texttt{0x2412}?}

\textbf{Procedimiento.} Decodificar \texttt{0x2412} usando los $9$
bits altos para \texttt{vpn} y $11$ bits bajos para \texttt{offset}.

\quad \texttt{0x2412} = $0000\ 0010\ 0100\ 0001\ 0010_2$ (20 bits)

\quad \texttt{vpn} (9 bits altos) = $0\ 0000\ 0100_2$ = $4$

\quad \texttt{offset} (11 bits bajos) = $100\ 0001\ 0010_2$ = \texttt{0x412} = $1042$

Alternativa numérica:

\quad \texttt{0x2412} $\div$ \texttt{0x800} = $9234 \div 2048 = 4$ (cociente)

\quad \texttt{0x2412} $\bmod$ \texttt{0x800} = $9234 - 4 \times 2048 = 1042$ = \texttt{0x412}

\textbf{Consulta a la tabla.} La página $\texttt{vpn} = 4$ está en el
rango $[3, 253]$ del \emph{gap}, donde todas las entradas tienen
$V = 0$.

\begin{respuesta}
\textbf{(3.IV)} La dirección \texttt{0x2412} \textbf{no se mapea a
ninguna página física}: cae en \texttt{vpn} $= 4$, que pertenece al
\textbf{gap inválido} entre los datos globales y el stack. El acceso
generaría un \emph{page fault} (V=0) y el kernel enviaría
\texttt{SIGSEGV} al proceso. No corresponde a código, datos ni stack.

\textbf{Nota didáctica}: el ejercicio testea que el estudiante
\emph{verifique} si la página es válida antes de asumir el área. Si la
página fuese de $4$\,KB (no $2$\,KB), la dirección \texttt{0x2412}
caería en \texttt{vpn} = $2$, que es la página de datos; pero con
$2$\,KB cae en el gap.
\end{respuesta}

\subsection*{3.V --- Página de guarda para detectar stack overflow}

\textbf{Idea.} Si el SO permite que el stack crezca dinámicamente
hacia direcciones más bajas (asignando bajo demanda las páginas $253$,
$252$, \ldots), conviene reservar una página \textbf{inmediatamente
después} de los datos globales como \emph{guarda}: nunca se mapea
($V=0$), y cualquier intento de tocarla provoca un page fault que el
kernel detecta como \emph{stack overflow}.

\textbf{Elección concreta.} La página de guarda más natural es
\texttt{vpn} $= 3$, justo después de los datos (\texttt{vpn} $= 2$):

\begin{itemize}[leftmargin=*,nosep]
    \item Marcar \texttt{vpn} $= 3$ con \textbf{$V = 0$ permanentemente}.
    \item Cuando el stack baja a \texttt{vpn} $= 4$, todavía hay
        espacio.
    \item Cuando intenta acceder a \texttt{vpn} $= 3$, el hardware
        levanta un fault, el kernel reconoce que la dirección está
        ``por debajo'' del límite válido del stack y \emph{aborta} el
        proceso con \emph{stack overflow} \textbf{antes} de pisar los
        datos.
\end{itemize}

\textbf{Espacio de direcciones lógicas que caen en la guarda.}
\texttt{vpn} $= 3$ con páginas de $2$\,KB cubre el rango:

\quad $[3 \times 2048, 4 \times 2048) = [6144, 8192) = $
\texttt{[0x1800, 0x2000)}

\begin{respuesta}
\textbf{(3.V)} La entrada \texttt{vpn} $= 3$ debe quedar sin mapear
($V = 0$) y servir como guarda. El espacio de direcciones lógicas
inválidas asociado es:

\begin{center}
\Large \texttt{[0x1800, 0x2000)}
\end{center}

\noindent es decir, los $2$\,KB que ocuparía esa página entre el área
de datos y la zona donde podría crecer el stack.
\end{respuesta}

% =====================================================================
\section{Ejercicio 4 --- Sistema de archivos FAT32 (2 pts)}
% =====================================================================

\textbf{Enunciado.} Completar el diagrama del sistema de archivos
FAT32 con bloques de datos de $512$ bytes, correspondiente al árbol:

\begin{lstlisting}[style=textstyle]
\
|-- shell.com    (programa, 2 bloques de datos)
+-- docs         (directorio vacio)
\end{lstlisting}

\noindent Cada entrada de directorio ocupa $32$ bytes.

\subsection*{Análisis previo: cuántos bloques ocupa cada cosa}

\textbf{Tamaño de los datos de cada entrada del árbol.}

\begin{itemize}[leftmargin=*]
    \item \texttt{\textbackslash} (root directory): contiene $4$
        entradas (\texttt{.}, \texttt{..}, \texttt{shell.com},
        \texttt{docs}). $4 \times 32 = 128$ bytes. Cabe en
        \textbf{$1$ bloque} de $512$.
    \item \texttt{shell.com}: $750$ bytes según la tabla de root.
        $\lceil 750 / 512 \rceil = 2$ bloques.
    \item \texttt{docs}: directorio vacío significa que contiene
        sólo \texttt{.} y \texttt{..}: $2 \times 32 = 64$ bytes.
        Cabe en \textbf{$1$ bloque}.
\end{itemize}

Total ocupado: $1 + 2 + 1 = 4$ bloques (de $5$ disponibles).

\subsection*{Asignación propuesta de clusters}

\begin{center}
\small
\begin{tabular}{@{}cl@{}}
\toprule
\# de cluster & Contenido \\
\midrule
$1$ & directorio raíz (\texttt{\textbackslash}) \\
$2$ & \texttt{shell.com}, bloque 1 (bytes $0$--$511$) \\
$3$ & directorio \texttt{docs} \\
$4$ & \texttt{shell.com}, bloque 2 (bytes $512$--$749$ y \emph{padding}) \\
$5$ & libre \\
\bottomrule
\end{tabular}
\end{center}

\subsection*{Tabla FAT completada}

Cada entrada de la FAT dice qué cluster sigue al actual, o si es el
final de la cadena (\texttt{EOC}) o si está libre ($0$).

\begin{center}
\begin{tikzpicture}[every node/.style={font=\ttfamily\small},
    cell/.style={draw,minimum width=1.4cm,minimum height=0.8cm,align=center}]
\node[cell,fill=c3!30] (f1) at (0,0)    {EOC};
\node[cell,fill=c2!30,right=0pt of f1]  (f2) {$4$};
\node[cell,fill=c4!30,right=0pt of f2]  (f3) {EOC};
\node[cell,fill=c2!30,right=0pt of f3]  (f4) {EOC};
\node[cell,fill=cfree!30,right=0pt of f4] (f5) {$0$};

\node[font=\scriptsize] at (f1) [yshift=-0.55cm] {1};
\node[font=\scriptsize] at (f2) [yshift=-0.55cm] {2};
\node[font=\scriptsize] at (f3) [yshift=-0.55cm] {3};
\node[font=\scriptsize] at (f4) [yshift=-0.55cm] {4};
\node[font=\scriptsize] at (f5) [yshift=-0.55cm] {5};

\node[font=\bfseries] at ($(f1.north west)+(-0.7,0)$) {FAT};
\end{tikzpicture}
\end{center}

\textbf{Justificación entrada por entrada:}

\begin{itemize}[leftmargin=*,nosep]
    \item \texttt{FAT[1] = EOC}: el directorio raíz ocupa exactamente
        un bloque, así que la cadena termina en cluster $1$.
    \item \texttt{FAT[2] = 4}: \texttt{shell.com} comienza en
        cluster $2$ y \emph{encadena} al cluster $4$ para su segundo
        bloque (notar que $3$ está ocupado por \texttt{docs}).
    \item \texttt{FAT[3] = EOC}: \texttt{docs} ocupa un solo bloque.
    \item \texttt{FAT[4] = EOC}: cluster $4$ es el último bloque de
        \texttt{shell.com}.
    \item \texttt{FAT[5] = 0}: cluster libre.
\end{itemize}

\subsection*{Contenido de los bloques de datos}

\textbf{Cluster 1 (directorio raíz \texttt{\textbackslash}):}

\begin{center}
\small
\begin{tabular}{@{}llc@{}}
\toprule
nombre & tipo / tamaño & 1.\textsuperscript{er} cluster \\
\midrule
\texttt{.}          & dir, size = 128 & $1$ \\
\texttt{..}         & dir, size = 128 & $1$ \\
\texttt{shell.com}  & file, size = 750 & $2$ \\
\texttt{docs}       & dir, size = 64  & $3$ \\
\bottomrule
\end{tabular}
\end{center}

Notar que \texttt{..} de la raíz apunta a sí misma (cluster $1$),
porque la raíz no tiene padre.

\textbf{Cluster 2:} primeros $512$ bytes del programa \texttt{shell.com}.

\textbf{Cluster 3 (directorio \texttt{docs}):}

\begin{center}
\small
\begin{tabular}{@{}llc@{}}
\toprule
nombre & tipo / tamaño & 1.\textsuperscript{er} cluster \\
\midrule
\texttt{.}  & dir, size = 64  & $3$ \\
\texttt{..} & dir, size = 128 & $1$ \\
\bottomrule
\end{tabular}
\end{center}

\noindent\texttt{..} de \texttt{docs} apunta a la raíz (cluster $1$);
\texttt{.} apunta al propio cluster del directorio.

\textbf{Cluster 4:} los $238$ bytes finales de \texttt{shell.com}
(bytes $512$--$749$) más \emph{padding} hasta completar los $512$
bytes del cluster.

\textbf{Cluster 5:} libre, sin contenido válido.

\subsection*{Diagrama final del disco}

\begin{center}
\begin{tikzpicture}[every node/.style={font=\ttfamily\small},
    cell/.style={draw,minimum width=1.9cm,minimum height=1.1cm,align=center}]
\node[font=\bfseries] at (-1.2,0.5) {data:};
\node[cell,fill=c3!30] (d1) at (0,0)    {raíz \\ \scriptsize(\texttt{\textbackslash})};
\node[cell,fill=c2!30,right=0pt of d1]  (d2) {shell.com \\ \scriptsize bloque 1};
\node[cell,fill=c4!30,right=0pt of d2]  (d3) {docs};
\node[cell,fill=c2!30,right=0pt of d3]  (d4) {shell.com \\ \scriptsize bloque 2};
\node[cell,fill=cfree!30,right=0pt of d4] (d5) {libre};
\node[font=\scriptsize] at (d1.south) [yshift=-0.18cm] {1};
\node[font=\scriptsize] at (d2.south) [yshift=-0.18cm] {2};
\node[font=\scriptsize] at (d3.south) [yshift=-0.18cm] {3};
\node[font=\scriptsize] at (d4.south) [yshift=-0.18cm] {4};
\node[font=\scriptsize] at (d5.south) [yshift=-0.18cm] {5};

\draw[->, c2!70!black,thick] (d2.north) to[bend left=30] (d4.north);
\end{tikzpicture}
\end{center}

La flecha arriba muestra el encadenamiento de \texttt{shell.com}:
cluster $2$ $\to$ cluster $4$ (saltando cluster $3$ que está ocupado
por \texttt{docs}). Este ``salto'' es justamente la utilidad de la
FAT: la tabla permite fragmentar archivos a lo largo del disco sin
que el cliente se entere.

\begin{respuesta}
La tabla FAT queda \textbf{[EOC, 4, EOC, EOC, 0]} para los clusters
$1$ a $5$. Los bloques de datos contienen, en orden: directorio raíz,
primer bloque de \texttt{shell.com}, directorio \texttt{docs}, segundo
bloque de \texttt{shell.com}, libre.
\end{respuesta}

\begin{idea}
La elección de poner \texttt{docs} en el cluster $3$ (en lugar de
poner \texttt{shell.com} consecutivo en clusters $2$--$3$) ilustra
cómo la fragmentación natural del filesystem se gestiona con la FAT.
En un disco real con miles de archivos creándose y borrándose, esta
estructura permite ubicar bloques en cualquier lado y reconstruir la
cadena en tiempo $O(\text{longitud})$ con la tabla.
\end{idea}

% =====================================================================
\section{Ejercicio 5 --- V/F sobre usuarios y permisos (1.5 pts)}
% =====================================================================

\textbf{Enunciado.} Determinar la verdad o falsedad de las siguientes
sentencias. Justificar las falsas.

\subsection*{5.a --- ``Un proceso se ejecuta con los privilegios del
usuario real''}

\begin{respuesta}
\textbf{F}. Un proceso se ejecuta con los privilegios del
\textbf{usuario efectivo} (\texttt{euid}), no del real
(\texttt{ruid}). En un proceso ``normal'' los dos coinciden, pero el
modelo UNIX deliberadamente los separa para soportar programas con
\textbf{bit SUID}: cuando se ejecuta un binario con \texttt{u+s}, el
kernel pone como \texttt{euid} del proceso el dueño del archivo
ejecutable, no el del usuario que lo lanzó.

Ejemplo clásico: \texttt{/usr/bin/passwd} pertenece a root con el bit
SUID; cuando un usuario común lo ejecuta, su \texttt{ruid} sigue
siendo el suyo pero su \texttt{euid} pasa a ser root, lo que le
permite escribir en \texttt{/etc/shadow}. Los chequeos de permisos
del kernel se hacen contra \texttt{euid}, no contra \texttt{ruid}.
\end{respuesta}

\subsection*{5.b --- ``\texttt{sudo} solicita la contraseña del
usuario ejecutor''}

\emph{``El comando \texttt{sudo} permite ejecutar un programa con
los privilegios de otro usuario y solicita la contraseña del usuario
ejecutor.''}

\begin{respuesta}
\textbf{V}. Las dos partes son correctas:

\begin{itemize}[leftmargin=*,nosep]
    \item \texttt{sudo} permite ejecutar un comando como otro usuario
        (por default \texttt{root}; con \texttt{-u} se elige otro
        target).
    \item Por defecto, \texttt{sudo} pide la \textbf{contraseña del
        usuario que lo invoca} (\emph{the user typing sudo}), no la
        del usuario objetivo. La verificación es ``demostrame que sos
        vos quien está pidiendo el privilegio''; el archivo
        \texttt{/etc/sudoers} decide a quién está autorizado a
        elevarse.
\end{itemize}

\textbf{Nota}: existe la opción \texttt{Defaults targetpw} que
invierte el comportamiento (pide la del target), pero no es el
default.
\end{respuesta}

\subsection*{5.c --- ``Es posible recuperar la contraseña\ldots''}

\emph{``Es posible recuperar la contraseña de un usuario ya que el
sistema la conoce.''}

\begin{respuesta}
\textbf{F}. El sistema \textbf{no conoce la contraseña en texto
plano}. Lo que guarda en \texttt{/etc/shadow} es un \textbf{hash}
criptográfico (típicamente \texttt{bcrypt}, \texttt{argon2},
\texttt{sha512crypt}) con \emph{salt}. La función hash es de un solo
sentido: dado el hash es computacionalmente inviable recuperar la
contraseña original.

Cuando el usuario hace login, el sistema toma la contraseña que
ingresa, le aplica la misma función hash con el mismo salt, y
compara con el hash guardado. Si coinciden, autentica.

Por eso, si un usuario olvida su contraseña, el administrador no
puede ``ver cuál era''; solo puede \textbf{asignarle una nueva}
(\texttt{passwd <usuario>}). Esto es una propiedad de seguridad
deliberada: aún si un atacante roba \texttt{/etc/shadow}, no
obtiene contraseñas en claro.
\end{respuesta}

\subsection*{5.d --- Permisos \texttt{rw-rw-r-{}-}}

\emph{``Un archivo con permisos \texttt{rw-rw-r-{}-} puede ser escrito
por el dueño y cualquier usuario del grupo asignado al archivo pero
no por el resto de usuarios.''}

\begin{respuesta}
\textbf{V}. Decodifiquemos el modo \texttt{rw-rw-r-{}-}:

\begin{center}
\small
\begin{tabular}{@{}llll@{}}
\toprule
Categoría & Permisos & Octal & Significado \\
\midrule
Owner (\texttt{u})  & \texttt{rw-} & $6$ & leer y escribir \\
Group (\texttt{g})  & \texttt{rw-} & $6$ & leer y escribir \\
Other (\texttt{o})  & \texttt{r-{}-} & $4$ & solo leer \\
\bottomrule
\end{tabular}
\end{center}

Total: $664$ en octal. Por tanto:

\begin{itemize}[leftmargin=*,nosep]
    \item El \textbf{dueño} puede escribir ($w$ en \texttt{u}).
    \item Los miembros del \textbf{grupo asignado} al archivo pueden
        escribir ($w$ en \texttt{g}).
    \item Los \textbf{demás usuarios} (no dueño ni del grupo) solo
        pueden leer.
\end{itemize}

La afirmación coincide exactamente. \textbf{V.}
\end{respuesta}

% =====================================================================
\section*{Resumen del parcial}
\addcontentsline{toc}{section}{Resumen del parcial}
% =====================================================================

\begin{center}
\small
\begin{tabular}{@{}clp{8.5cm}@{}}
\toprule
\# & Tema & Respuesta sintética \\
\midrule
1 & V/F memoria        & a:V, b:V, c:V, d:F (\texttt{mmap} es lazy), e:V \\
2 & Buddy system        & ptr1@0 (512), ptr3@640 (128), ptr4@1024 (1024); libres: 128@512 y 256@768 \\
3 & Paginación          & 512 ptes, vpn 9 bits + offset 11 bits; tabla con código 0--1, datos 2, stack 254--255, gap 3--253; \texttt{0x2412} cae en gap (page fault); guarda en vpn 3 con rango \texttt{[0x1800, 0x2000)} \\
4 & FAT32               & FAT = [EOC, 4, EOC, EOC, 0]; raíz @1, shell.com @2 $\to$ 4, docs @3 \\
5 & V/F usuarios        & a:F (\texttt{euid}, no \texttt{ruid}); b:V; c:F (hash); d:V \\
\bottomrule
\end{tabular}
\end{center}

\end{document}
